\documentclass{article}
\usepackage[utf8]{inputenc}
\usepackage{geometry}
\usepackage{enumitem}
\usepackage{hyperref}
\usepackage{xcolor}
\usepackage{graphicx}
\usepackage{array}
\usepackage{multirow}
\usepackage{amsmath, amssymb}
\usepackage{chemfig}
\usepackage{mhchem}
\usepackage{tikz}
\usepackage{setspace}
\usepackage{soul}
\usepackage{mathtools}
\usepackage{booktabs}
\usepackage{chemformula}
\usepackage{siunitx}
\usetikzlibrary{positioning, calc}
\usepackage{longtable}
\geometry{margin=1in}
\setlength{\parskip}{0.5em}
\usepackage{cancel}


\begin{document}

\begin{center}
    \vspace*{-1cm}
    \begin{tabular}{p{12cm} r}
        & \textbf{Name:} \underline{\hspace{5cm}} \\
    \end{tabular}

    \vspace{0.2cm}
    {\LARGE \textbf{Week 9 Worksheet}}
\end{center}    
\\


\noindent\begin{tabular}{|p{0.9\linewidth}|}
\hline
\vspace{0.2cm}
\textbf{(1)} How many moles are in 25.0~g of NaCl?\\[0.2cm]

\begin{minipage}[t]{0.85\linewidth}
\begin{enumerate}[label=\alph*.]
\item 0.427
\item 2.50
\item 0.125
\item 1.00
\end{enumerate}
\end{minipage}\\[0.2cm]
\hline
\end{tabular}

\vspace{0.2cm}
\noindent\begin{tabular}{|p{0.9\linewidth}|}
\hline
\textbf{(2)} $2\text{H}_2 + \text{O}_2 \rightarrow 2\text{H}_2\text{O}$. If 3.0~mol of H$_2$ react with excess O$_2$, how many moles of H$_2$O are formed?\\[0.2cm]

\begin{minipage}[t]{0.85\linewidth}
\begin{enumerate}[label=\alph*.]
\item 1.5
\item 3.0
\item 6.0
\item 2.0
\end{enumerate}
\end{minipage}\\[0.2cm]
\hline
\end{tabular}

\vspace{0.2cm}
\noindent\begin{tabular}{|p{0.9\linewidth}|}
\hline
\textbf{(3)} $\text{N}_2 + 3\text{H}_2 \rightarrow 2\text{NH}_3$. If 4.0~mol H$_2$ react with 2.0~mol N$_2$, which is the limiting reagent?\\[0.2cm]

\begin{minipage}[t]{0.85\linewidth}
\begin{enumerate}[label=\alph*.]
\item N$_2$
\item H$_2$
\item Neither
\item Both
\end{enumerate}
\end{minipage}\\[0.2cm]
\hline
\end{tabular}

\vspace{0.2cm}
\noindent\begin{tabular}{|p{0.9\linewidth}|}
\hline
\textbf{(4)} In the reaction $\text{C} + \text{O}_2 \rightarrow \text{CO}_2$, how many grams of CO$_2$ form when 10.0~g of carbon reacts completely?\\[0.2cm]

\begin{minipage}[t]{0.85\linewidth}
\begin{enumerate}[label=\alph*.]
\item 22.0
\item 36.7
\item 10.0
\item 44.0
\end{enumerate}
\end{minipage}\\[0.2cm]
\hline
\end{tabular}
\newpage

\noindent\begin{tabular}{|p{0.9\linewidth}|}
\hline
\textbf{(5)} $\text{Ca} + 2\text{H}_2\text{O} \rightarrow \text{Ca(OH)}_2 + \text{H}_2$. How many liters of H$_2$ are produced at 298~K and 1.50~atm from 2.00~mol Ca?\\[0.2cm]

\begin{minipage}[t]{0.85\linewidth}
\begin{enumerate}[label=\alph*.]
\item 36.5
\item 44.8
\item 24.7
\item 32.6
\end{enumerate}
\end{minipage}\\[0.2cm]
\hline
\end{tabular}

\vspace{0.2cm}
\noindent\begin{tabular}{|p{0.9\linewidth}|}
\hline
\textbf{(6)} $2\text{K} + \text{Cl}_2 \rightarrow 2\text{KCl}$. How many moles of KCl are formed from 0.500~mol Cl$_2$?\\[0.2cm]

\begin{minipage}[t]{0.85\linewidth}
\begin{enumerate}[label=\alph*.]
\item 0.250
\item 1.00
\item 0.500
\item 2.00
\end{enumerate}
\end{minipage}\\[0.2cm]
\hline
\end{tabular}

\vspace{0.2cm}
\noindent\begin{tabular}{|p{0.9\linewidth}|}
\hline
\textbf{(7)} How many molecules are present in 0.300~mol of CO$_2$?\\[0.2cm]

\begin{minipage}[t]{0.85\linewidth}
\begin{enumerate}[label=\alph*.]
\item $1.81\times10^{23}$
\item $3.01\times10^{23}$
\item $6.02\times10^{23}$
\item $9.03\times10^{23}$
\end{enumerate}
\end{minipage}\\[0.2cm]
\hline
\end{tabular}

\vspace{0.2cm}
\noindent\begin{tabular}{|p{0.9\linewidth}|}
\hline
\textbf{(8)} If 50.0~g of Fe reacts with excess O$_2$ to form Fe$_2$O$_3$, and the actual yield is 68.0~g, what is the percent yield if the theoretical yield was 80.0~g?\\[0.2cm]

\begin{minipage}[t]{0.85\linewidth}
\begin{enumerate}[label=\alph*.]
\item 85.0\%
\item 90.0\%
\item 68.0\%
\item 80.0\%
\end{enumerate}
\end{minipage}\\[0.2cm]
\hline
\end{tabular}

\newpage

\noindent\begin{tabular}{|p{0.9\linewidth}|}
\hline
\textbf{(9)} $\text{Zn} + 2\text{HCl} \rightarrow \text{ZnCl}_2 + \text{H}_2$. How many grams of ZnCl$_2$ form when 0.500~mol of Zn reacts completely?\\[0.2cm]

\begin{minipage}[t]{0.85\linewidth}
\begin{enumerate}[label=\alph*.]
\item 34.0
\item 68.0
\item 136
\item 17.0
\end{enumerate}
\end{minipage}\\[0.2cm]
\hline
\end{tabular}

\vspace{0.2cm}
\noindent\begin{tabular}{|p{0.9\linewidth}|}
\hline
\textbf{(10)} How many oxygen atoms are in 0.200~mol of Al$_2$O$_3$?\\[0.2cm]

\begin{minipage}[t]{0.85\linewidth}
\begin{enumerate}[label=\alph*.]
\item $3.61\times10^{23}$
\item $1.20\times10^{23}$
\item $6.02\times10^{23}$
\item $7.23\times10^{23}$
\end{enumerate}
\end{minipage}\\[0.2cm]
\hline
\end{tabular}

\vspace{0.2cm}
\noindent\begin{tabular}{|p{0.9\linewidth}|}
\hline
\textbf{(11)} $2\text{H}_2 + \text{O}_2 \rightarrow 2\text{H}_2\text{O}$. If 4.0~mol H$_2$ react with 2.0~mol O$_2$, which is limiting?\\[0.2cm]

\begin{minipage}[t]{0.85\linewidth}
\begin{enumerate}[label=\alph*.]
\item H$_2$
\item O$_2$
\item Neither
\item Both
\end{enumerate}
\end{minipage}\\[0.2cm]
\hline
\end{tabular}

\vspace{0.2cm}
\noindent\begin{tabular}{|p{0.9\linewidth}|}
\hline
\textbf{(12)} If the theoretical yield of NH$_3$ from $\text{N}_2 + 3\text{H}_2 \rightarrow 2\text{NH}_3$ is 10.0~g, and the actual yield is 7.80~g, what is the percent yield?\\[0.2cm]

\begin{minipage}[t]{0.85\linewidth}
\begin{enumerate}[label=\alph*.]
\item 78.0\%
\item 100\%
\item 82.0\%
\item 70.0\%
\end{enumerate}
\end{minipage}\\[0.2cm]
\hline
\end{tabular}

\newpage

\noindent\begin{tabular}{|p{0.9\linewidth}|}
\hline
\textbf{(13)} How many moles of oxygen gas are needed to completely burn 2.0~mol CH$_4$? ($\text{CH}_4 + 2\text{O}_2 \rightarrow \text{CO}_2 + 2\text{H}_2\text{O}$)\\[0.2cm]

\begin{minipage}[t]{0.85\linewidth}
\begin{enumerate}[label=\alph*.]
\item 1.0
\item 2.0
\item 4.0
\item 8.0
\end{enumerate}
\end{minipage}\\[0.2cm]
\hline
\end{tabular}

\vspace{0.2cm}
\noindent\begin{tabular}{|p{0.9\linewidth}|}
\hline
\textbf{(14)} $2\text{Na} + \text{Cl}_2 \rightarrow 2\text{NaCl}$. If 5.00~mol Na react with 1.50~mol Cl$_2$, what is the limiting reagent?\\[0.2cm]

\begin{minipage}[t]{0.85\linewidth}
\begin{enumerate}[label=\alph*.]
\item Na
\item Cl$_2$
\item Neither
\item Both
\end{enumerate}
\end{minipage}\\[0.2cm]
\hline
\end{tabular}

\vspace{0.2cm}
\noindent\begin{tabular}{|p{0.9\linewidth}|}
\hline
\textbf{(15)} If 0.750.~mol CaCO$_3$ decomposes to $\text{CaO} + \text{CO}_2$, how many moles of CO$_2$ are produced?\\[0.2cm]

\begin{minipage}[t]{0.85\linewidth}
\begin{enumerate}[label=\alph*.]
\item 1.50
\item 0.750
\item 0.375
\item 0.250
\end{enumerate}
\end{minipage}\\[0.2cm]
\hline
\end{tabular}

\vspace{0.2cm}
\noindent\begin{tabular}{|p{0.9\linewidth}|}
\hline
\textbf{(16)} $\text{Zn} + \text{S} \rightarrow \text{ZnS}$. If 3.00~mol Zn react with 2.00~mol S, which is the limiting reagent?\\[0.2cm]

\begin{minipage}[t]{0.85\linewidth}
\begin{enumerate}[label=\alph*.]
\item Zn
\item S
\item Neither
\item Both
\end{enumerate}
\end{minipage}\\[0.2cm]
\hline
\end{tabular}

\newpage

\noindent\begin{tabular}{|p{0.9\linewidth}|}
\hline
\textbf{(17)} If 5.00~mol of O$_2$ react with excess H$_2$, how many grams of H$_2$O are formed? ($2\text{H}_2 + \text{O}_2 \rightarrow 2\text{H}_2\text{O}$)\\[0.2cm]

\begin{minipage}[t]{0.85\linewidth}
\begin{enumerate}[label=\alph*.]
\item 90.0
\item 45.0
\item 180
\item 22.4
\end{enumerate}
\end{minipage}\\[0.2cm]
\hline
\end{tabular}

\vspace{0.2cm}
\noindent\begin{tabular}{|p{0.9\linewidth}|}
\hline
\textbf{(18)} $2\text{Al} + 3\text{Cl}_2 \rightarrow 2\text{AlCl}_3$. If 2.0~mol Al react with 4.0~mol Cl$_2$, what is the limiting reagent?\\[0.2cm]

\begin{minipage}[t]{0.85\linewidth}
\begin{enumerate}[label=\alph*.]
\item Al
\item Cl$_2$
\item Neither
\item Both
\end{enumerate}
\end{minipage}\\[0.2cm]
\hline
\end{tabular}

\vspace{0.2cm}
\noindent\begin{tabular}{|p{0.9\linewidth}|}
\hline
\textbf{(19)} $2\text{H}_2\text{O}_2 \rightarrow 2\text{H}_2\text{O} + \text{O}_2$. If 6.00~mol H$_2$O$_2$ decompose, how many moles of O$_2$ form?\\[0.2cm]

\begin{minipage}[t]{0.85\linewidth}
\begin{enumerate}[label=\alph*.]
\item 3.00
\item 6.00
\item 12.0
\item 1.50
\end{enumerate}
\end{minipage}\\[0.2cm]
\hline
\end{tabular}

\vspace{0.2cm}
\noindent\begin{tabular}{|p{0.9\linewidth}|}
\hline
\textbf{(20)} What is the percent yield if 10.0~g of product are expected but only 8.20~g are obtained?\\[0.2cm]

\begin{minipage}[t]{0.85\linewidth}
\begin{enumerate}[label=\alph*.]
\item 82.0\%
\item 92.0\%
\item 80.0\%
\item 88.0\%
\end{enumerate}
\end{minipage}\\[0.2cm]
\hline
\end{tabular}

\newpage

\noindent\begin{tabular}{|p{0.9\linewidth}|}
\hline
\textbf{(21)} $2\text{HCl} + \text{Mg} \rightarrow \text{MgCl}_2 + \text{H}_2$. How many moles of H$_2$ are produced from 5.00~mol of HCl?\\[0.2cm]

\begin{minipage}[t]{0.85\linewidth}
\begin{enumerate}[label=\alph*.]
\item 5.00
\item 2.50
\item 10.0
\item 1.00
\end{enumerate}
\end{minipage}\\[0.2cm]
\hline
\end{tabular}

\vspace{0.2cm}
\noindent\begin{tabular}{|p{0.9\linewidth}|}
\hline
\textbf{(22)} Complete the combustion reaction of methane and calculate how many grams of CO$_2$ are produced from 0.500~mol CH$_4$: \\[0.2cm]
\hspace{1em} $\text{CH}_4 + \text{O}_2 \rightarrow \text{CO}_2 + \text{H}_2\text{O}$ \\[0.2cm]

\begin{minipage}[t]{0.85\linewidth}
\begin{enumerate}[label=\alph*.]
\item 11.0
\item 22.0
\item 44.0
\item 16.0
\end{enumerate}
\end{minipage}\\[0.2cm]
\hline
\end{tabular}

\vspace{0.2cm}
\noindent\begin{tabular}{|p{0.9\linewidth}|}
\hline
\textbf{(23)} $2\text{Na} + 2\text{H}_2\text{O} \rightarrow 2\text{NaOH} + \text{H}_2$. If 4.00~mol Na react, how many moles of H$_2$ are released?\\[0.2cm]

\begin{minipage}[t]{0.85\linewidth}
\begin{enumerate}[label=\alph*.]
\item 1.00
\item 2.00
\item 4.00
\item 8.00
\end{enumerate}
\end{minipage}\\[0.2cm]
\hline
\end{tabular}

\vspace{0.2cm}
\noindent\begin{tabular}{|p{0.9\linewidth}|}
\hline
\textbf{(24)} In the reaction $\text{CaCO}_3 \rightarrow \text{CaO} + \text{CO}_2$, if 50.0~g of CaCO$_3$ produce 25.0~g of CaO, what is the percent yield if 28.0~g was expected?\\[0.2cm]

\begin{minipage}[t]{0.85\linewidth}
\begin{enumerate}[label=\alph*.]
\item 89.3\%
\item 75.0\%
\item 92.0\%
\item 100\%
\end{enumerate}
\end{minipage}\\[0.2cm]
\hline
\end{tabular}

\newpage

\noindent\begin{tabular}{|p{0.9\linewidth}|}
\hline
\textbf{(25)} $\text{Fe}_2\text{O}_3 + 3\text{CO} \rightarrow 2\text{Fe} + 3\text{CO}_2$. If 2.00~mol Fe$_2$O$_3$ react with 6.00~mol CO, what mass of Fe is formed assuming 100\% yield?\\[0.2cm]

\begin{minipage}[t]{0.85\linewidth}
\begin{enumerate}[label=\alph*.]
\item 56.0
\item 112
\item 224
\item 168
\end{enumerate}
\end{minipage}\\[0.2cm]
\hline
\end{tabular}

\vspace{0.2cm}
\noindent\begin{tabular}{|p{0.9\linewidth}|}
\hline
\textbf{(26)} A 425~mL cylinder contains nitrogen at $2.50\times10^{2}$~kPa and 22.0$^\circ$C. What pressure (in atm) will result if the gas is moved to a 1.00~L container and heated to 125$^\circ$C, assuming moles constant?\\[0.2cm]

\begin{minipage}[t]{0.85\linewidth}
\begin{enumerate}[label=\alph*.]
\item 0.564~atm
\item 1.41~atm
\item 0.282~atm
\item 2.24~atm
\end{enumerate}
\end{minipage}\\[0.2cm]
\hline
\end{tabular}

\vspace{0.2cm}
\noindent\begin{tabular}{|p{0.9\linewidth}|}
\hline
\textbf{(27)} 0.375~mol of an ideal gas occupies $6.50\times10^{2}$~mL at 27.0$^\circ$C. Using the ideal gas law, what is the pressure in mmHg? \\[0.2cm]

\begin{minipage}[t]{0.85\linewidth}
\begin{enumerate}[label=\alph*.]
\item 575~mmHg
\item 760.~mmHg
\item 430.~mmHg
\item $1.08\times10^{4}$~mmHg
\end{enumerate}
\end{minipage}\\[0.2cm]
\hline
\end{tabular}

\vspace{0.2cm}
\noindent\begin{tabular}{|p{0.9\linewidth}|}
\hline
\textbf{(28)} A sealed syringe has 8.00~mL of gas at 295~K and 750.~torr. If the plunger is pulled to 24.0~mL and temperature falls to 270.~K, what is the new pressure (in kPa)?\\[0.2cm]

\begin{minipage}[t]{0.85\linewidth}
\begin{enumerate}[label=\alph*.]
\item 125~kPa
\item 188~kPa
\item 62.5~kPa
\item 30.5~kPa
\end{enumerate}
\end{minipage}\\[0.2cm]
\hline
\end{tabular}

\newpage

\noindent\begin{tabular}{|p{0.9\linewidth}|}
\hline
\textbf{(29)} 1.50~g of helium is confined in a 2.00~L flask at 350.~K. Using the ideal gas law, what is the pressure in atm? \\[0.2cm]

\begin{minipage}[t]{0.85\linewidth}
\begin{enumerate}[label=\alph*.]
\item 0.650.~atm
\item 5.39~atm
\item 0.325~atm
\item 1.45~atm
\end{enumerate}
\end{minipage}\\[0.2cm]
\hline
\end{tabular}

\vspace{0.2cm}
\noindent\begin{tabular}{|p{0.9\linewidth}|}
\hline
\textbf{(30)} A gas at 1.20~atm and 300~K fills a 5.00~L vessel. If the temperature is increased to 450.~K and pressure is allowed to equilibrate at constant moles, what is the final volume (L)?\\[0.2cm]

\begin{minipage}[t]{0.85\linewidth}
\begin{enumerate}[label=\alph*.]
\item 7.50~L
\item 3.33~L
\item 5.00~L
\item 11.3~L
\end{enumerate}
\end{minipage}\\[0.2cm]
\hline
\end{tabular}

\vspace{0.2cm}
\noindent\begin{tabular}{|p{0.9\linewidth}|}
\hline
\textbf{(31)} A sample of gas has a volume of 2.50~L at a pressure of 1.20~atm. If the pressure is increased to 3.00~atm at constant temperature, what is the new volume of the gas? (Boyle’s Law — indirect relationship)\\[0.2cm]

\begin{minipage}[t]{0.85\linewidth}
\begin{enumerate}[label=\alph*.]
\item 1.00~L
\item 2.10.~L
\item 0.950~L
\item 3.75~L
\end{enumerate}
\end{minipage}\\[0.2cm]
\hline
\end{tabular}

\vspace{0.2cm}
\noindent\begin{tabular}{|p{0.9\linewidth}|}
\hline
\textbf{(32)} A gas mixture in a 10.0~L tank exerts total pressure 4.00~atm. The partial pressures are $P_A=1.50$~atm and $P_B=0.75$~atm. What is $P_C$ and how many moles of C are present at 25.0$^\circ$C?\\[0.2cm]

\begin{minipage}[t]{0.85\linewidth}
\begin{enumerate}[label=\alph*.]
\item $P_C=1.75$~atm, $n_C=7.22\times10^{-1}$~mol
\item $P_C=1.75$~atm, $n_C=1.63$~mol
\item $P_C=1.25$~atm, $n_C=0.722$~mol
\item $P_C=1.75$~atm, $n_C=3.61\times10^{-1}$~mol
\end{enumerate}
\end{minipage}\\[0.2cm]
\hline
\end{tabular}

\newpage

\noindent\begin{tabular}{|p{0.9\linewidth}|}
\hline
\textbf{(33)} A gas initially at 640.~mmHg and 20.0$^\circ$C is compressed to 300~mL from 750.~mL while temperature changes to 60.0$^\circ$C. What is final pressure in mmHg?\\[0.2cm]

\begin{minipage}[t]{0.85\linewidth}
\begin{enumerate}[label=\alph*.]
\item $1.12\times10^{3}$~mmHg
\item 512~mmHg
\item 290.~mmHg
\item $1.50\times10^{3}$~mmHg
\end{enumerate}
\end{minipage}\\[0.2cm]
\hline
\end{tabular}

\vspace{0.2cm}
\noindent\begin{tabular}{|p{0.9\linewidth}|}
\hline
\textbf{(34)} 0.250~mol of oxygen gas is collected at 740.~mmHg and 18.0$^\circ$C in a 2.00~L flask. What would be the pressure at 0.500~mol in same flask and same T? (assume ideal)\\[0.2cm]

\begin{minipage}[t]{0.85\linewidth}
\begin{enumerate}[label=\alph*.]
\item 1480.~mmHg
\item 740.~mmHg
\item 370.~mmHg
\item 1110.~mmHg
\end{enumerate}
\end{minipage}\\[0.2cm]
\hline
\end{tabular}

\vspace{0.2cm}
\noindent\begin{tabular}{|p{0.9\linewidth}|}
\hline
\textbf{(35)} A gas at 1.00~atm and 298~K is allowed to expand isothermally from 2.00~L to 6.00~L. What is final pressure in atm?\\[0.2cm]

\begin{minipage}[t]{0.85\linewidth}
\begin{enumerate}[label=\alph*.]
\item 3.00~atm
\item 0.333~atm
\item 0.667~atm
\item 1.00~atm
\end{enumerate}
\end{minipage}\\[0.2cm]
\hline
\end{tabular}

\vspace{0.2cm}
\noindent\begin{tabular}{|p{0.9\linewidth}|}
\hline
\textbf{(36)} A sample of gas weighing 4.80~g exerts 0.900~atm in a 2.50~L container at 300.~K. What is the molar mass of the gas? \\[0.2cm]

\begin{minipage}[t]{0.85\linewidth}
\begin{enumerate}[label=\alph*.]
\item 58.2~g~mol$^{-1}$
\item 120.~g~mol$^{-1}$
\item 52.5~g~mol$^{-1}$
\item 16.0~g~mol$^{-1}$
\end{enumerate}
\end{minipage}\\[0.2cm]
\hline
\end{tabular}

\newpage

\noindent\begin{tabular}{|p{0.9\linewidth}|}
\hline
\textbf{(37)} At 1.00~atm and 298~K, a gas sample has density 0.860~g/L. What is its molar mass?\\[0.2cm]

\begin{minipage}[t]{0.85\linewidth}
\begin{enumerate}[label=\alph*.]
\item 19.4~g~mol$^{-1}$
\item 24.7~g~mol$^{-1}$
\item 29.0~g~mol$^{-1}$
\item 256~g~mol$^{-1}$
\end{enumerate}
\end{minipage}\\[0.2cm]
\hline
\end{tabular}

\vspace{0.2cm}
\noindent\begin{tabular}{|p{0.9\linewidth}|}
\hline
\textbf{(38)} A 150.~mL sample of gas at 25$^\circ$C and 755~mmHg contains $2.00\times10^{21}$ molecules. What is Avogadro's number implied by these measurements? (students must convert units and compute $N_A$)\\[0.2cm]

\begin{minipage}[t]{0.85\linewidth}
\begin{enumerate}[label=\alph*.]
\item $6.02\times10^{23}$
\item $5.12\times10^{23}$
\item $7.30\times10^{23}$
\item $3.01\times10^{23}$
\end{enumerate}
\end{minipage}\\[0.2cm]
\hline
\end{tabular}

\vspace{0.2cm}
\noindent\begin{tabular}{|p{0.9\linewidth}|}
\hline
\textbf{(39)} A gas occupies 500.~mL at 25$^\circ$C. What volume will it occupy at 100.$^\circ$C if the pressure remains constant? \\[0.2cm]

\begin{minipage}[t]{0.85\linewidth}
\begin{enumerate}[label=\alph*.]
\item 583~mL
\item 625~mL
\item 665~mL
\item 740.~mL
\end{enumerate}
\end{minipage}\\[0.2cm]
\hline
\end{tabular}

\vspace{0.2cm}
\noindent\begin{tabular}{|p{0.9\linewidth}|}
\hline
\textbf{(40)} A gas initially at 1.20~atm, 298~K and 250.~mL is heated to 350.~K and allowed to expand until pressure is 1.00~atm. What is the final volume (mL)?\\[0.2cm]

\begin{minipage}[t]{0.85\linewidth}
\begin{enumerate}[label=\alph*.]
\item 330.~mL
\item 250.~mL
\item 290.~mL
\item 400~mL
\end{enumerate}
\end{minipage}\\[0.2cm]
\hline
\end{tabular}

\newpage

\noindent\begin{tabular}{|p{0.9\linewidth}|}
\hline
\textbf{(41)} A gas is stored in a container at 2.00~atm and 20.0$^\circ$C. If the temperature is increased to 80.0$^\circ$C at constant volume, what will be the new pressure? \\[0.2cm]

\begin{minipage}[t]{0.85\linewidth}
\begin{enumerate}[label=\alph*.]
\item 2.35~atm
\item 2.60.~atm
\item 2.80.~atm
\item 3.00~atm
\end{enumerate}
\end{minipage}\\[0.2cm]
\hline
\end{tabular}

\vspace{0.2cm}
\noindent\begin{tabular}{|p{0.9\linewidth}|}
\hline
\textbf{(42)} A cylinder contains 0.850~mol of gas at 2.50~atm. What is the temperature (K) if the volume is 15.0~L?\\[0.2cm]

\begin{minipage}[t]{0.85\linewidth}
\begin{enumerate}[label=\alph*.]
\item 210.~K
\item 298~K
\item 425~K
\item 512~K
\end{enumerate}
\end{minipage}\\[0.2cm]
\hline
\end{tabular}

\vspace{0.2cm}
\noindent\begin{tabular}{|p{0.9\linewidth}|}
\hline
\textbf{(43)} A sealed vessel at 1.00~atm and 300.~K contains $3.00\times10^{22}$ molecules. If vessel is evacuated and recharged with the same number of molecules at 600.~K, what will be the new pressure (atm) in the same vessel?\\[0.2cm]

\begin{minipage}[t]{0.85\linewidth}
\begin{enumerate}[label=\alph*.]
\item 0.500~atm
\item 1.00~atm
\item 2.00~atm
\item 4.00~atm
\end{enumerate}
\end{minipage}\\[0.2cm]
\hline
\end{tabular}

\vspace{0.2cm}
\noindent\begin{tabular}{|p{0.9\linewidth}|}
\hline
\textbf{(44)} True or False: Gas particles are in constant, random motion.\\[0.2cm]

\begin{minipage}[t]{0.85\linewidth}
\begin{enumerate}[label=\alph*.]
\item True
\item False
\end{enumerate}
\end{minipage}\\[0.2cm]
\hline
\end{tabular}

\newpage

\noindent\begin{tabular}{|p{0.9\linewidth}|}
\hline
\textbf{(45)} True or False: Gas pressure is caused by collisions of particles with the container walls.\\[0.2cm]

\begin{minipage}[t]{0.85\linewidth}
\begin{enumerate}[label=\alph*.]
\item True
\item False
\end{enumerate}
\end{minipage}\\[0.2cm]
\hline
\end{tabular}

\vspace{0.2cm}
\noindent\begin{tabular}{|p{0.9\linewidth}|}
\hline
\textbf{(46)} Increasing the temperature of a gas increases the average:\\[0.2cm]

\begin{minipage}[t]{0.85\linewidth}
\begin{enumerate}[label=\alph*.]
\item Pressure only
\item Molecular speed
\item Volume only
\item Mass
\end{enumerate}
\end{minipage}\\[0.2cm]
\hline
\end{tabular}

\vspace{0.2cm}
\noindent\begin{tabular}{|p{0.9\linewidth}|}
\hline
\textbf{(47)} True or False: Lighter gas particles move faster than heavier ones at the same temperature.\\[0.2cm]

\begin{minipage}[t]{0.85\linewidth}
\begin{enumerate}[label=\alph*.]
\item True
\item False
\end{enumerate}
\end{minipage}\\[0.2cm]
\hline
\end{tabular}

\vspace{0.2cm}
\noindent\begin{tabular}{|p{0.9\linewidth}|}
\hline
\textbf{(48)} Gas particles collide with each other:\\[0.2cm]

\begin{minipage}[t]{0.85\linewidth}
\begin{enumerate}[label=\alph*.]
\item Inelastically, losing energy each time
\item Perfectly elastically, no net energy loss
\item Only with container walls
\item Only at high pressure
\end{enumerate}
\end{minipage}\\[0.2cm]
\hline
\end{tabular}

\newpage

\noindent\begin{tabular}{|p{0.9\linewidth}|}
\hline
\textbf{(49)} True or False: Real gases behave more like ideal gases at low pressure.\\[0.2cm]

\begin{minipage}[t]{0.85\linewidth}
\begin{enumerate}[label=\alph*.]
\item True
\item False
\end{enumerate}
\end{minipage}\\[0.2cm]
\hline
\end{tabular}

\vspace{0.2cm}
\noindent\begin{tabular}{|p{0.9\linewidth}|}
\hline
\textbf{(50)} A gas sample occupies 1.20~L when it contains 0.050~mol of gas. What will its volume be if 0.200~mol of gas is added at constant temperature and pressure? \\[0.2cm]

\begin{minipage}[t]{0.85\linewidth}
\begin{enumerate}[label=\alph*.]
\item 2.40.~L
\item 3.60.~L
\item 4.80.~L
\item 6.00~L
\end{enumerate}
\end{minipage}\\[0.2cm]
\hline
\end{tabular}

\vspace{0.2cm}
\noindent\begin{tabular}{|p{0.9\linewidth}|}
\hline
\textbf{(51)} True or False: Gas particles are so small that their volume is negligible compared to the volume of the container.\\[0.2cm]

\begin{minipage}[t]{0.85\linewidth}
\begin{enumerate}[label=\alph*.]
\item True
\item False
\end{enumerate}
\end{minipage}\\[0.2cm]
\hline
\end{tabular}

\vspace{0.2cm}
\noindent\begin{tabular}{|p{0.9\linewidth}|}
\hline
\textbf{(52)} True or False: All gases at the same temperature have the same average kinetic energy.\\[0.2cm]

\begin{minipage}[t]{0.85\linewidth}
\begin{enumerate}[label=\alph*.]
\item True
\item False
\end{enumerate}
\end{minipage}\\[0.2cm]
\hline
\end{tabular}

\newpage

\noindent\begin{tabular}{|p{0.9\linewidth}|}
\hline
\textbf{(53)} Increasing the number of gas particles in a container at constant volume and temperature:\\[0.2cm]

\begin{minipage}[t]{0.85\linewidth}
\begin{enumerate}[label=\alph*.]
\item Increases pressure
\item Decreases molecular speed
\item Lowers temperature
\item Decreases collisions
\end{enumerate}
\end{minipage}\\[0.2cm]
\hline
\end{tabular}

\vspace{0.2cm}
\noindent\begin{tabular}{|p{0.9\linewidth}|}
\hline
\textbf{(54)} True or False: Increasing volume at constant temperature decreases gas pressure.\\[0.2cm]

\begin{minipage}[t]{0.85\linewidth}
\begin{enumerate}[label=\alph*.]
\item True
\item False
\end{enumerate}
\end{minipage}\\[0.2cm]
\hline
\end{tabular}

\vspace{0.2cm}
\noindent\begin{tabular}{|p{0.9\linewidth}|}
\hline
\textbf{(55)} True or False: Gas particles attract each other significantly in an ideal gas.\\[0.2cm]

\begin{minipage}[t]{0.85\linewidth}
\begin{enumerate}[label=\alph*.]
\item True
\item False
\end{enumerate}
\end{minipage}\\[0.2cm]
\hline
\end{tabular}

\vspace{0.2cm}
\noindent\begin{tabular}{|p{0.9\linewidth}|}
\hline
\textbf{(56)} True or False: The faster a gas particle moves, the more kinetic energy it has.\\[0.2cm]

\begin{minipage}[t]{0.85\linewidth}
\begin{enumerate}[label=\alph*.]
\item True
\item False
\end{enumerate}
\end{minipage}\\[0.2cm]
\hline
\end{tabular}

\newpage

\noindent\begin{tabular}{|p{0.9\linewidth}|}
\hline
\textbf{(57)} Which of the following increases when temperature increases for a gas at constant volume?\\[0.2cm]

\begin{minipage}[t]{0.85\linewidth}
\begin{enumerate}[label=\alph*.]
\item Pressure
\item Molecular mass
\item Intermolecular forces
\item Volume
\end{enumerate}
\end{minipage}\\[0.2cm]
\hline
\end{tabular}

\vspace{0.2cm}
\noindent\begin{tabular}{|p{0.9\linewidth}|}
\hline
\textbf{(58)} A gas has a volume of 3.00~L at 760.~mmHg and 300.~K. What will be its volume at 600.~mmHg and 350.~K? \\[0.2cm]

\begin{minipage}[t]{0.85\linewidth}
\begin{enumerate}[label=\alph*.]
\item 4.50.~L
\item 5.00~L
\item 5.25~L
\item 6.00~L
\end{enumerate}
\end{minipage}\\[0.2cm]
\hline
\end{tabular}

\vspace{0.2cm}
\noindent\begin{tabular}{|p{0.9\linewidth}|}
\hline
\textbf{(59)} True or False: Gas particle collisions are perfectly elastic, meaning no kinetic energy is lost.\\[0.2cm]

\begin{minipage}[t]{0.85\linewidth}
\begin{enumerate}[label=\alph*.]
\item True
\item False
\end{enumerate}
\end{minipage}\\[0.2cm]
\hline
\end{tabular}

\vspace{0.2cm}
\noindent\begin{tabular}{|p{0.9\linewidth}|}
\hline
\textbf{(60)} True or False: Gas particles are very far apart compared to their size.\\[0.2cm]

\begin{minipage}[t]{0.85\linewidth}
\begin{enumerate}[label=\alph*.]
\item True
\item False
\end{enumerate}
\end{minipage}\\[0.2cm]
\hline
\end{tabular}

\end{document}
